% ----------------------------------------------------------------------------
\chapter{Coarea, slicing and traces}\label{ch:coarea}
\module{NoCompromise.BV.AnnularGluingFull,
        NoCompromise.BV.BoundaryTraces,
        NoCompromise.BV.Coarea,
        NoCompromise.BV.CoareaCritical,
        NoCompromise.BV.CoareaL1Lebesgue,
        NoCompromise.BV.CoareaRegular,
        NoCompromise.BV.ExactCuts,
        NoCompromise.BV.Gluing,
        NoCompromise.BV.GoodRadii,
        NoCompromise.BV.GoodTruncation,
        NoCompromise.BV.JumpDisintegration,
        NoCompromise.BV.Slicing,
        NoCompromise.BV.WeightedCoarea}

\section{Coarea for BV}\label{sec:coarea}

\begin{theorem}[Weighted $C^1$ coarea]\label{thm:weighted-coarea}
Let $n\ge2$, let $u\in C^1(U)$ with $U\subset\R^n$ open, and let
$g:U\to[0,\infty]$ be Borel.
Then
\begin{equation}\label{eq:weighted-coarea}
  \int_Ug(x)\,|\nabla u(x)|\,dx
  =\int_\R\int_{U\cap\{u=t\}}g\,d\Hs^{n-1}\,dt .
\end{equation}
No Sard theorem is used; critical levels are permitted.
\end{theorem}

\begin{proof}\uses{lem:coarea-regular-part,lem:coarea-critical-part}
Combine Lemmas~\ref{lem:coarea-regular-part} and
\ref{lem:coarea-critical-part}: the first proves
\eqref{eq:weighted-coarea} with $U\setminus Z$ in place of $U$, where
$Z:=\{\nabla u=0\}$; the second shows $Z$ contributes zero to both sides.
\end{proof}

\begin{lemma}[The regular part]\label{lem:coarea-regular-part}
Let $n\ge2$, let $u\in C^1(U)$ with $U\subset\R^n$ open, let
$g:U\to[0,\infty]$ be Borel, and put $Z:=\{\nabla u=0\}$.  Then
\[
  \int_{U\setminus Z}g\,|\nabla u|\,dx
  =\int_\R\int_{(U\setminus Z)\cap\{u=t\}}g\,d\Hs^{n-1}\,dt .
\]
\end{lemma}

\begin{proof}
Cover $U\setminus Z$ by countably many relatively compact boxes $V$ on each
of which one fixed derivative, after permuting coordinates say
$\partial_nu$, has constant nonzero sign, subdividing so that each line
parallel to $e_n$ meets $V$ in an interval on which $u$ is strictly
monotone.  Writing $x'=(x_1,\ldots,x_{n-1})$, on such a box
$\Psi(x',x_n):=(x',u(x))$ is a $C^1$ diffeomorphism onto its image with
Jacobian $|\partial_nu|$, and \ref{base:calculus} gives
\begin{equation}\label{eq:coarea-cov}
  \int_Vg|\nabla u|
  =\int_{\Psi(V)}g(\Psi^{-1}(y',t))
     \frac{|\nabla u|}{|\partial_nu|}(\Psi^{-1}(y',t))\,dy'\,dt .
\end{equation}
For fixed $t$ the inverse graph $x_n=\gamma_t(x')$ satisfies
\begin{equation}\label{eq:coarea-graph}
  \nabla\gamma_t=-\frac{(\partial_1u,\ldots,\partial_{n-1}u)}{\partial_nu},
  \qquad
  \sqrt{1+|\nabla\gamma_t|^2}=\frac{|\nabla u|}{|\partial_nu|},
\end{equation}
so the inner integral in \eqref{eq:coarea-cov} is the graph definition of
$\int_{V\cap\{u=t\}}g\,d\Hs^{n-1}$.  Disjointify the cover, apply the identity
first to simple $g$, and use monotone convergence.
\end{proof}

\begin{lemma}[The critical set contributes nothing]\label{lem:coarea-critical-part}
Let $n\ge2$, let $u\in C^1(U)$ with $U\subset\R^n$ open, and put
$Z:=\{\nabla u=0\}$.  Then
\[
  \Hs^{n-1}\bigl(Z\cap\{u=t\}\bigr)=0
  \qquad\text{for a.e.\ }t\in\R .
\]
\end{lemma}

\begin{proof}
Fix $K\Subset U$; if $K\cap Z=\varnothing$ there is nothing to prove.
Consider the grid cubes of side $\delta$ meeting $K\cap Z$, with $\delta$
small enough that these lie in $U$, and put
\[
  \omega_K(\delta):=\sup\{|\nabla u(x)|:\dist(x,K\cap Z)\le\sqrt n\,\delta\}.
\]
Uniform continuity of $\nabla u$ near $K$ and $\nabla u=0$ on $Z$ give
$\omega_K(\delta)\to0$.  The $u$-image of each selected cube has length at
most $\sqrt n\,\delta\,\omega_K(\delta)$.  For fixed $t$, the selected cubes
whose images contain $t$ cover $K\cap Z\cap\{u=t\}$, so their
$(n-1)$-st powers of the diameters bound its $(n-1)$-dimensional Hausdorff
content at scale $\sqrt n\delta$.  Integrating in $t$,
\begin{equation}\label{eq:coarea-content}
  \int_\R\Hs^{n-1}_{\sqrt n\delta}
       \bigl(K\cap Z\cap\{u=t\}\bigr)\,dt
  \le C_n\omega_K(\delta)
       \!\!\sum_{Q\cap K\cap Z\ne\varnothing}\!\!\delta^n
  \le C_K\,\omega_K(\delta).
\end{equation}
As $\delta\downarrow0$ Hausdorff contents increase to $\Hs^{n-1}$ while the
right side tends to zero; monotone convergence and an exhaustion of $Z$
finish.
\end{proof}

\begin{corollary}[Coarea with an $L^1$ weight]\label{cor:coarea-L1}
\uses{thm:weighted-coarea}
Let $n\ge2$, let $u\in C^1(U)$ with $U\subset\R^n$ open, let $A\subset U$
be measurable, and let $G\in L^1(A)$.  Then
\begin{equation}\label{eq:coarea-L1}
  \int_{A\cap\{|\nabla u|>0\}}G\,dx
  =\int_\R\int_{A\cap\{u=t\}\cap\{|\nabla u|>0\}}
     \frac{G}{|\nabla u|}\,d\Hs^{n-1}\,dt .
\end{equation}
If $|\nabla u|>0$ throughout $A$ both restrictions may be omitted.  This is
the form used in Chapter~\ref{ch:appK}.
\end{corollary}

\begin{proof}\uses{thm:weighted-coarea}
Extend $G$ by zero off $A$ and apply Theorem~\ref{thm:weighted-coarea} to
$g:=|G|/|\nabla u|$ on $A\cap\{|\nabla u|>0\}$ and $g:=0$ elsewhere; then
split $G$ into positive and negative parts.
\end{proof}

\begin{theorem}[Coarea for BV]\label{thm:bv-coarea}
\uses{def:variation,def:perimeter}
Let $n\ge2$, let $U\subset\R^n$ be open, and let
$f\in BV_{\mathrm{loc}}(U)$.  Then
\begin{equation}\label{eq:bv-coarea}
  |Df|(U)=\int_\R\Per_n(\{f>t\};U)\,dt .
\end{equation}
\end{theorem}

\begin{proof}
\uses{lem:strict-approx,lem:lsc,thm:weighted-coarea}
\emph{Smooth case.} Theorem~\ref{thm:weighted-coarea} with $g=1$ gives
\eqref{eq:bv-coarea} for $f\in C^1$.

\emph{$\le$ for BV.} If $|Df|(U)=\infty$, the asserted upper bound for the
level-set integral is automatic.  Otherwise take $f_j$ as in
Theorem~\ref{lem:strict-approx}.  On each member of a countable relatively
compact exhaustion, layer cake and Fubini turn $f_j\to f$ in $L^1$ there
into, after one diagonal subsequence,
$\ind{\{f_j>t\}}\to\ind{\{f>t\}}$ in $\Lloc$ for a.e.\ $t$.  Lemma~\ref{lem:lsc},
Fatou and \eqref{eq:strict-approx} give
$\int_\R\Per_n(\{f>t\};U)\,dt\le|Df|(U)$.

\emph{$\ge$ for BV.} First let $0\le f\le M$.  For $X\in C^1_c(U;\R^n)$,
Fubini and the layer-cake formula give
\[
  \int_Uf\,\dvg X=\int_0^M\!\!\int_{\{f>t\}}\dvg X\,dx\,dt
  \le\|X\|_\infty\int_0^M\Per_n(\{f>t\};U)\,dt ;
\]
take the supremum over $\|X\|_\infty\le1$.  A bounded signed function is
reduced to this case by adding a constant, whose integral against
$\dvg X$ is zero, and shifting the level parameter.

For a general $f$, put $T_M(f):=\max\{-M,\min\{f,M\}\}$.  The truncation
chain rule
\begin{equation}\label{eq:bv-truncation-chain}
  |D(T_Mf)|(A)\le |Df|(A)\qquad(A\subset U\text{ open})
\end{equation}
follows first on every $V\Subset A$ by applying smooth $1$-Lipschitz
approximations of $T_M$ to a local strict smooth approximation of $f$, using
the ordinary chain rule and Lemma~\ref{lem:lsc}, and then by exhausting $A$.
Since $T_Mf\to f$ in $L^1_{\rm loc}$,
Lemma~\ref{lem:lsc} and \eqref{eq:bv-truncation-chain} give
$|D(T_Mf)|(U)\to|Df|(U)$ (with extended values allowed).  Apply the bounded
case to $T_Mf$ and let $M\to\infty$; Tonelli and monotone convergence in the
level variable give the reverse inequality for $f$.
\end{proof}

\begin{corollary}[Spherical and planar slicing]\label{cor:slicing}
\uses{thm:weighted-coarea,def:density}
Let $E$ have locally finite perimeter.  For $x_0\in\R^3$ and
$0\le a<b<\infty$,
\[
  \int_a^b\Hs^2\bigl(E^{(1)}\cap\partial B_r(x_0)\bigr)\,dr
  =\bigl|E\cap(B_b(x_0)\setminus B_a(x_0))\bigr| ,
\]
and, for $\nu\in\Sph^2$ and $a<b$ in $\R$,
\[
 \int_a^b\Hs^2\bigl(E^{(1)}\cap\{x:x\cdot\nu=t\}\bigr)\,dt
 =\bigl|E\cap\{x:a<x\cdot\nu<b\}\bigr|.
\]
\end{corollary}

\begin{proof}\uses{thm:weighted-coarea,def:density}
The sets $E$ and $E^{(1)}$ agree up to a Lebesgue-null set.  Apply
Theorem~\ref{thm:weighted-coarea} with weight $g=\ind{E^{(1)}}$, first with
$u(x)=|x-x_0|$ off the null point $x_0$ and then with $u(x)=x\cdot\nu$.
In both cases $|\nabla u|=1$ on the relevant set, and restriction to the
indicated interval gives the two identities.
\end{proof}

\begin{lemma}[One-dimensional BV slices]\label{lem:bv-slices}
\uses{def:variation}
Let $E\subset\R^n$ have locally finite perimeter, $n\ge2$.  For a.e.\ line $\ell$
parallel to a fixed direction, $t\mapsto\ind E(x'+t e)$ is a
one-dimensional BV function; the number of its jumps is finite on compact
intervals for a.e.\ $\ell$, and
$D_e\ind E$ disintegrates as the one-dimensional jump measures over $\ell$.
\end{lemma}

\begin{proof}\uses{thm:polar,lem:strict-approx}
Apply Fubini to \eqref{eq:variation} with test fields of the form
$\phi(x')\psi(t)e$; strict approximation
(Theorem~\ref{lem:strict-approx}) identifies the resulting slice variations
with the disintegration of $|D_e\ind E|$.
\end{proof}

\section{Traces, cut identities and gluing}\label{sec:traces}

\begin{theorem}[Traces on a $C^1$ domain]\label{thm:traces}
\uses{def:variation,def:density}
Let $A\subset\R^3$ be a bounded $C^1$ domain and
$f\in BV_{\mathrm{loc}}(\R^3)$.  Then $f$ has interior and exterior traces
$\Tr_Af,\Tr_{A^c}f\in L^1(\partial A)$, and
\begin{equation}\label{eq:trace-product}
  D(\ind Af)=\ind A\,Df-\Tr_Af\,\nu_A\,\Hs^2\llcorner\partial A .
\end{equation}
It also satisfies the exterior formula
\begin{equation}\label{eq:trace-product-exterior}
  D(\ind{\R^3\setminus\overline A}f)
  =\ind{\R^3\setminus\overline A}\,Df
    +\Tr_{A^c}f\,\nu_A\,\Hs^2\llcorner\partial A .
\end{equation}
\end{theorem}

\begin{proof}\uses{lem:bv-slices,lem:strict-approx,lem:bilip-chain}
The bi-Lipschitz BV chain rule (Lemma~\ref{lem:bilip-chain}, whose proof uses
only strict approximation and change of variables) puts $f$ in $BV$ after
flattening a $C^1$ boundary chart.  Fubini and one-dimensional BV slicing
(Lemma~\ref{lem:bv-slices}) then give one-sided limits on a.e.\ normal line;
changing between overlapping charts preserves these limits by the same chain
rule and \ref{base:calculus}.
For \eqref{eq:trace-product} and \eqref{eq:trace-product-exterior}, integrate
by parts first for smooth $f$ on the two sides and then along a strict
approximation (Theorem~\ref{lem:strict-approx}).  There is no missing boundary
term in the limit: in a graph chart it is exactly the endpoint term in the
one-dimensional fundamental theorem of calculus, and
$L^1(\partial A)$ convergence of the one-sided traces follows by integrating
that one-dimensional identity.
\end{proof}

\begin{definition}[Good radius]\label{def:good-radius}
\uses{def:density,def:perimeter}
Let $E$ have locally finite perimeter and $x_0\in\R^3$.  Call $r>0$ a
\emph{good radius} for $E$ at $x_0$ if
\[
  |D\ind E|(\partial B_r(x_0))=0
  \quad\text{and}\quad
  \Tr_{B_r(x_0)}\ind E=\ind{E^{(1)}}\ \ \Hs^2\text{-a.e.\ on }\partial B_r(x_0).
\]
\end{definition}

\begin{lemma}[Almost every radius is good]\label{lem:good-radius-ae}
\uses{def:good-radius}
For $E$ of locally finite perimeter and any $x_0$, a.e.\ $r>0$ is a good
radius for $E$ at $x_0$.
\end{lemma}

\begin{proof}\uses{cor:slicing,thm:traces}
The set of $r$ with $|D\ind E|(\partial B_r(x_0))>0$ is countable, since
$|D\ind E|$ is locally finite.  The trace identity holds for a.e.\ $r$ by
Corollary~\ref{cor:slicing} together with Theorem~\ref{thm:traces} and
Lebesgue's density theorem.
\end{proof}

\begin{proposition}[Exact cut identities]\label{prop:cut-identities}
\uses{def:good-radius}
Let $E$ have locally finite perimeter and let $r$ be a good radius for $E$ at
$x_0$.  Then
\begin{align}
  \Per(E\cap B_r(x_0))
    &=\Per(E;B_r(x_0))+\Hs^2\bigl(E^{(1)}\cap\partial B_r(x_0)\bigr),
      \label{eq:cut-in}\\
  \Per(E\setminus B_r(x_0))
    &=\Per\bigl(E;\R^3\setminus\overline{B_r(x_0)}\bigr)
      +\Hs^2\bigl(E^{(1)}\cap\partial B_r(x_0)\bigr).
      \label{eq:cut-out}
\end{align}
The corresponding identities hold for halfspaces $\{x\cdot\nu<\ell\}$ at
a.e.\ $\ell$.
\end{proposition}

\begin{proof}\uses{thm:traces,lem:good-radius-ae}
Apply \eqref{eq:trace-product} with $A=B_r(x_0)$.  At a good radius the two
measures on its right side are mutually singular, since
$|D\ind E|(\partial B_r)=0$; taking total variations gives \eqref{eq:cut-in},
and \eqref{eq:trace-product-exterior} gives \eqref{eq:cut-out}; no
boundedness of the exterior domain is being assumed.  For halfspaces use
$x\mapsto x\cdot\nu$ in place of
$x\mapsto|x-x_0|$.
\end{proof}

\begin{fnote}[Why exactness matters]\label{fnote:exact-cut}
Only the \emph{equalities} \eqref{eq:cut-in}--\eqref{eq:cut-out}, not
inequalities, make the perimeter error of a cut equal to twice the section
area.  Weakening them to inequalities loses a sign in
Theorem~\ref{thm:decomposition} and in Lemma~\ref{lem:two-sided-density}.
\end{fnote}

\begin{lemma}[Good truncation]\label{lem:good-truncation}
\uses{def:good-radius,def:perimeter}
Let $E\subset\R^3$ have finite measure and finite perimeter.  There are radii
$R_n\uparrow\infty$, each a good radius for $E$ at the origin, such that
\begin{equation}\label{eq:good-truncation}
  a_n:=\Hs^2\bigl(E^{(1)}\cap\partial B_{R_n}\bigr)\longrightarrow0,
  \qquad
  |E\setminus B_{R_n}|\longrightarrow0,
  \qquad
  \Per(E\cap B_{R_n})\longrightarrow\Per(E).
\end{equation}
\end{lemma}

\begin{proof}\uses{cor:slicing,lem:good-radius-ae,prop:cut-identities,lem:lsc}
Corollary~\ref{cor:slicing} gives
\[
  \int_n^{n+1}\Hs^2\bigl(E^{(1)}\cap\partial B_r\bigr)\,dr
  =\bigl|E\cap(B_{n+1}\setminus B_n)\bigr|\longrightarrow0 ,
\]
so by Lemma~\ref{lem:good-radius-ae} choose $R_n\in(n,n+1)$ among the good
radii with $a_n\to0$; then $|E\setminus B_{R_n}|\to0$ because $|E|<\infty$.
Since $R_n$ is strictly increasing, $\ind{E\cap B_{R_n}}\uparrow\ind E$ a.e.
The trace identity \eqref{eq:cut-in} gives
$\Per(E\cap B_{R_n})=\Per(E;B_{R_n})+a_n$, hence
$\limsup_n\Per(E\cap B_{R_n})\le\Per(E)$, while global $L^1$ convergence and
Lemma~\ref{lem:lsc} give the reverse $\liminf$.
\end{proof}

\begin{fnote}[Shared by two consumers]\label{fnote:good-truncation-shared}
Lemma~\ref{lem:good-truncation} is the common first step of
Lemma~\ref{lem:bounded-approx} and Theorem~\ref{thm:smooth-approx}.  It is
extracted here so that neither proof re-derives it, and so that the
truncation step carries no dependence on the Coulomb energy.
\end{fnote}

\begin{proposition}[Gluing formula]\label{prop:gluing}
\uses{thm:traces}
Let $U\subset\R^3$ be open, let $A\Subset U$ be a bounded $C^1$ domain, and
let $E,F$ have locally finite perimeter in $U$.  Put
$H:=(F\cap A)\cup(E\cap(U\setminus A))$.  Then
\begin{equation}\label{eq:gluing}
\begin{aligned}
  \Per(H;U)
  \le{}&\Per(F;A)+\Per\bigl(E;U\setminus\overline A\bigr)\\
  &+\int_{\partial A}|\Tr_AF-\Tr_{A^c}E|\,d\Hs^2 .
\end{aligned}
\end{equation}
For globally finite-perimeter $E,F$, taking $U=\R^3$ gives the global
formula.
\end{proposition}

\begin{proof}\uses{thm:traces}
Apply \eqref{eq:trace-product} on both sides of $\partial A$ and add the
distributional derivatives; the boundary measures combine into the trace
mismatch term.
\end{proof}

\begin{proposition}[Annular gluing at a good radius]\label{prop:annular-gluing}
\uses{prop:gluing,def:good-radius}
Let $A_{a,b}:=B_b\setminus\overline{B_a}$, suppose $E_j\to E$ in
$L^1(A_{a,b})$, suppose $E_j,E,F$ have locally finite perimeter, and for
a.e.\ $r\in(a,b)$ put
$H_{j,r}:=(F\cap B_r)\cup(E_j\setminus B_r)$.  Then
for every open $W$ with $\overline{B_b}\Subset W$,
\begin{equation}\label{eq:annular-gluing}
  \Per(H_{j,r};W)\le\Per(F;B_r)
   +\Per\bigl(E_j;W\setminus\overline{B_r}\bigr)
   +\Hs^2\bigl((F^{(1)}\triangle E_j^{(1)})\cap\partial B_r\bigr),
\end{equation}
and
\begin{equation}\label{eq:annular-coarea}
  \int_a^b\Hs^2\bigl((F^{(1)}\triangle E_j^{(1)})\cap\partial B_r\bigr)\,dr
  =\bigl|(F\triangle E_j)\cap A_{a,b}\bigr| .
\end{equation}
Consequently, if $F=E$ a.e. on the annulus, there are radii
$r_j\in(a,b)$ for which the last term of
\eqref{eq:annular-gluing} tends to zero; a diagonal choice handles expanding
annuli.  If all three sets have globally finite perimeter, exhaustion by
$W\uparrow\R^3$ gives the same formula with global perimeters.
\end{proposition}

\begin{proof}\uses{prop:gluing,cor:slicing,lem:good-radius-ae}
\eqref{eq:annular-gluing} is \eqref{eq:gluing} with $A=B_r$; the coarea
identity \eqref{eq:annular-coarea} is Corollary~\ref{cor:slicing} applied to
$F\triangle E_j$.  The averaging conclusion is Chebyshev in $r$.
\end{proof}

\begin{fnote}\label{fnote:no-gluing-theorem}
Proposition~\ref{prop:annular-gluing} is the precise content of every
``interpolate across a slicing annulus'' step in this document
(Lemma~\ref{lem:tangent-cone-minimizing}, Theorem~\ref{thm:decomposition},
Lemma~\ref{lem:height-bound}).  No separate gluing theorem may be invoked.
\end{fnote}

